% !PRINTED 179
\weberpage{179}

So geht beispielsweise jede lineare Function
\[
y = \overset{i}{\Sigma} a_i x_i
\]
in eine ebensolche Function
\[
y = \overset{i}{\Sigma} a_i' x_i'
\]
über, worin
\weq{3}{a_k' = \overset{i}{\Sigma} a_i \alpha_i^{(k)}.}
Betrachten wir $m$ solcher linearer Functionen
\[
y_\nu = \overset{i}{\Sigma} a_{i,\nu} x_i \qquad \nu = 1, 2 \ldots m,
\]
so erhalten wir $m$ transformirte Functionen
\[
y_\nu = \overset{i}{\Sigma} a_{i,\nu}' x_i' \qquad \nu = 1, 2 \ldots m.
\]
Die Determinanten der Systeme $y_\nu, y_\nu'$
\[
\Delta = \Sigma \pm a_{1,1} a_{2,2} \ldots a_{m,m}, \qquad \Delta' = \Sigma \pm a_{1,1}' a_{2,2}' \ldots a_{m,m}'
\]
stehen nach dem Multiplicationsgesetz der Determinanten in der Abhängigkeit
\weq{4}{\Delta' = r \Delta,}
so dass die eine nicht ohne die andere verschwindet.

\begin{center}
\S.\,56.

\gesperrt{Quadratische Formen.}
\end{center}

Von besonderer Wichtigkeit ist die Transformation der homogenen Functionen zweiten Grades oder der quadratischen Formen. Wir bezeichnen sie, wie schon im ersten Abschnitt, so:
\weq{1}{\varphi(x_1, x_2 \ldots x_m) = \overset{i,k}{\Sigma} a_{i,k} x_i x_k, \qquad a_{i,k} = a_{k,i}.}
Durch die Substitution \S.\,55, (1) geht $\varphi$ über in
\weq{2}{\Phi(x_1', x_2' \ldots x_m') = \overset{i,k}{\Sigma} a_{i,k}' x_i' x_k'.}
Setzen wir $x + \xi$ für $x$ und ebenso $x' + \xi'$ für $x'$, so dass die $\xi'$ mit den $\xi$ in derselben Abhängigkeit stehen, wie die $x'$ mit den $x$, so erhalten wir, wenn wir zur Abkürzung
\[
\tfrac{1}{2}\varphi'(x_i) = u_i, \qquad \tfrac{1}{2}\Phi'(x_i') = u_i'
\]
setzen, durch Vergleichung der Glieder der ersten Dimension (\S.\,16):
%CONT>>
